\subsection{Use case: a boundary time crystal under local dissipation}\label{sec:btc}

\paragraph{The question.} $N$ two-level systems with a collective drive and collective decay,
\begin{equation}\label{eq:btc}
  \dot\rho=-i[\omega J_x,\rho]+\frac{\kappa}{S}\,\mathcal D[J_-]\rho
  +\gamma\sum_{i=1}^N\mathcal D[\sigma_-^{(i)}]\rho+\gamma_\phi\sum_{i=1}^N\mathcal D[\sigma_z^{(i)}]\rho,\qquad
  \mathcal D[L]\rho=L\rho L^*-\tfrac12\{L^*L,\rho\},
\end{equation}
with $S=N/2$ and $J_a=\frac12\sum_i\sigma_a^{(i)}$. For $\gamma=\gamma_\phi=0$ this is the \emph{boundary time
crystal} of Iemini et al.~\cite{iemini2018}: for $\omega>\kappa$ the magnetisation $m_z=\langle J_z\rangle/S$
oscillates, the oscillations are damped at finite $N$ with a lifetime that grows linearly with
$N$~\cite{liu2026,nemeth2025}, and they persist as
$N\to\infty$ (a limit cycle of the mean-field equations, which become exact as $N\to\infty$ at fixed
times; Carollo and Lesanovsky~\cite{carollo2021} prove this for open Dicke models and note that the proof adapts to
all-to-all spin models and to individual spin dissipation). Real emitters also decay and dephase individually. Such
local terms do not conserve the total spin and, in the mean-field equations below, turn the limit cycle into a damped
oscillation, as Piccitto et al.~\cite{piccitto2021} found in mean field for jump operators acting collectively on
blocks of $N_s<N$ spins (damping rate $\propto1/N_s$, in a related $p,q$-spin model); robustness of
boundary time crystals has been studied in multi-level settings~\cite{prazeres2021}, and Tucker et
al.~\cite{tucker2018} studied single-particle dephasing in a related synchronised phase. We ask how long the
oscillations live as a function of $N$ and of the local rate, and how the two damping mechanisms---the collective
one, which vanishes as $N\to\infty$, and the local one, which does not---combine.

\paragraph{The algebra.} By Schur--Weyl duality, $(\C^2)^{\otimes N}=\bigoplus_jV_j\otimes\C^{m_j}$ with $\dim
V_j=2j+1$, $j=N/2,N/2-1,\dots\ge0$, and $m_j=\binom N{N/2-j}-\binom N{N/2-j-1}$, and the operators that commute with
the permutations of the sites form
\[
  \M=\bigoplus_jM_{2j+1}(\C)\otimes1_{m_j},\qquad \Tr_{(\C^2)^{\otimes N}}(A)=\sum_jm_j\Tr A_j .
\]
The right-hand side of~\eqref{eq:btc} commutes with the permutations, so a permutation-invariant state stays one:
$\rho=\bigoplus_j\sigma_j\otimes1_{m_j}/m_j$, where $\sigma=\bigoplus_j\sigma_j$ is the density with respect to
$\Tr_{\mathrm{blunt}}$ in the sense of Section~\ref{sec:math} (the ``Dicke-basis'' density of PIQS~\cite{shammah2018}).
The collective operators act inside the blocks, as the spin-$j$ matrices. A local term is a map between sectors: for
$L\in\{\sigma_-,\sigma_+,\sigma_z\}$ the map $\Phi(X)=\sum_iL^{(i)}XL^{(i)*}$ is normal and completely positive on
$\M$, and restricting to the last site and branching $N\to N-1$ ($\C^{m_j}$ is the direct sum of the multiplicity
spaces of $j_1=j\pm\frac12$ for $N-1$ sites, which $L^{(N)}$ preserves, and the $N$ sites contribute equally) gives
its Kraus operators on blunt densities,
\begin{equation}\label{eq:btc-kraus}
  \Phi(\sigma)_{j'}=\sum_j\sum_{j_1}K^{j'j}_{j_1}\sigma_jK^{j'j*}_{j_1},\qquad
  K^{j'j}_{j_1}=\Bigl(\frac{N\,m^{(N-1)}_{j_1}}{m^{(N)}_j}\Bigr)^{1/2}
  \bigl\langle j_1,\tfrac12;j'\bigr|\,1\otimes L\,\bigl|j_1,\tfrac12;j\bigr\rangle ,
\end{equation}
with $|j'-j|\le1$, $j_1\in\{j\pm\frac12\}\cap\{j'\pm\frac12\}$, and the coupled multiplets of $V_{j_1}\otimes\C^2$ in
the last factor; each $K$ has one non-zero diagonal, and $\Phi^*(1)=\sum_iL^{(i)*}L^{(i)}$ is $S+J_z$, $S-J_z$ or $N$.
(Closed forms of these coefficients are given in~\cite{chase2008,baragiola2010,shammah2018}; Eqs.~(40)--(46)
of~\cite{baragiola2010}, whose $\hat\sigma_z$ is half the Pauli matrix and whose $\rho_{J,M;J,M'}$ are the entries of
the blunt density, agree with~\eqref{eq:btc-kraus}. Permutational symmetry with individual decay and dephasing is also
exploited by the library PsiQuaSP~\cite{gegg2017} and by the method of Kirton and Keeling~\cite{kirton2017}.) In the library,
\texttt{permutational.PermutationalSystem(N)} is the algebra with these sectors; \texttt{local\_jump} returns $\Phi$ as an
\texttt{InterSectorChannel} of rank two, so \texttt{dynamics.lindblad\_evolve} handles~\eqref{eq:btc} with the
collective jumps as operators and the local ones between sectors, unchanged. Because every operator in~\eqref{eq:btc}
is banded in the basis $|j,m\rangle$, \texttt{PermutationalLindbladian} evaluates the same generator by shifted
elementwise products, $O(N^3)$ operations per state instead of $O(N^4)$ for block products, with all rates batched.
We drive about $y$ instead of $x$: the rotation $e^{-i\pi J_z/2}$ commutes with every dissipator, and the generator
becomes real, so the computation runs in the real algebra $\bigoplus_jM_{2j+1}(\mathbb R)$.
\IfFileExists{generated/permutational.tex}{For $N=\PMNmax$ a state has $\sum_j(2j+1)^2=\PMDimNmax$ entries, a density
matrix on $(\C^2)^{\otimes N}$ has $4^N\approx\PMDimFullNmax$.}{}

\noindent The program \repofile{permutational.py} makes one batched run over four local emission rates for $N=30$
in the real algebra $\bigoplus_jM_{2j+1}(\mathbb R)$, from all spins down, and returns $m_z(t)$ and the distribution
of $j$ (\texttt{PermutationalSystem}, \texttt{PermutationalLindbladian}, \texttt{coherent\_state},
\texttt{dynamics.rk4}, \texttt{m\_values}).

\paragraph{Validation.} \IfFileExists{generated/permutational.tex}{%
(i)~The banded generator, with all nine terms (drive about $x,y,z$, collective and local emission, pumping and
dephasing), equals the exact Lindblad generator on $(\C^2)^{\otimes N}$, compressed to $\M$ through Dicke bases built
independently from $J^2$ and $J_z$, to $\PMFullErr$ for $N\le\PMFullNmax$; an RK4 trajectory computed on
$(\C^2)^{\otimes\PMTrajN}$ gives the same blocks to $\PMTrajErr$.
(ii)~The generic route (\texttt{lindblad\_evolve} with \texttt{InterSectorChannel} jumps) agrees with the banded one
to $\PMLibErr$ ($N=\PMLibN$; $\PMLibSecGeneric$~s against $\PMLibSecBanded$~s).
(iii)~Against \texttt{qutip.piqs}~\cite{shammah2018} (QuTiP \PMPiqsVersion): the generators agree to a relative
$\PMPiqsGenErr$ (PIQS stores its rates in single precision, and its local dephasing rate is $4\gamma_\phi$); trajectories
of~\eqref{eq:btc} for $N=\PMPiqsNs$ agree to $\PMPiqsTrajErr$ in $m_z$ and to $\PMPiqsBlockErr$ in the blocks; the
steady states for $N=\PMPiqsSsNs$ to $\PMPiqsSsErr$.
(iv)~Pure collective decay $\mathcal D[J_-]$ (rate $\Gamma=1$) from $|S,S\rangle$ is a pure death process with rates
$\Gamma(S+m)(S-m+1)$, so $\int_0^\infty P(m>0)\,dt=\sum_{m=1}^S[\Gamma(S+m)(S-m+1)]^{-1}=H_N/(\Gamma(N+1))$ with
the harmonic number $H_N$; Table~\ref{tab:btc-sr} reproduces it to $\PMSRErr$ for $N\le\PMSRNmax$, together with the
superradiant delay $t_D\approx\ln N/(\Gamma N)$ and the reduction of the peak of the mean intensity below the
mean-field value $\Gamma N^2/4$ by the fluctuations of the delay~\cite{gross1982}.
(v)~Local emission, pumping and dephasing alone keep a product state a product of identical qubit states, so
$\langle J_z\rangle$, $\langle J_x\rangle$ and $\langle J^2\rangle$ follow from the Bloch equations of one qubit,
and the spin distribution is $p_j=m_j\sum_{m=-j}^j\lambda_+^{S+m}\lambda_-^{S-m}$ with the eigenvalues $\lambda_\pm$ of
the one-qubit state. For $N=\PMProdN$ the moments agree to a relative $\PMProdErr$ and $\sum_j|p_j-p_j^{\rm
exact}|=\PMProdPjErr$; this tests the multiplicities in~\eqref{eq:btc-kraus} far from the reach of the full space.
(vi)~Halving the time step changes $m_z$ by at most $\PMConvMz$ and $\tau$ by $\PMConvTau$ ($N=\PMConvN$).%
}{Run \texttt{scripts/permutational.py} to generate the numbers.}

\begin{table}[ht]
\centering\small
\caption{Pure collective decay from $|S,S\rangle$ ($\Gamma=1$): the integral of $P(m>0)$ against the exact value
$H_N/(N+1)$, the time $t_D$ of the peak of the mean intensity $\Gamma\langle J_+J_-\rangle$ in units of $\ln N/N$, and
the peak in units of $N^2/4$ (mean field: $1$ and $1$ for $N\to\infty$).}\label{tab:btc-sr}
\IfFileExists{generated/permutational.tex}{%
\begin{tabular}{rccccc}
\toprule
$N$ & $\int P(m>0)$ & $H_N/(N+1)$ & rel.\ dev. & $t_DN/\ln N$ & $I_{\max}/(N^2/4)$\\
\midrule
\PMSRRows
\bottomrule
\end{tabular}}{\fbox{\small generated/permutational.tex not found}}
\end{table}

\paragraph{The sweep.} We take $\omega=\IfFileExists{generated/permutational.tex}{\PMOmega}{2}\kappa$, all spins down
at $t=0$, and define the lifetime $\tau$ as the first time at which the half-swing of $m_z$ between consecutive
extrema falls below $1/e$ of the first one (log-linear interpolation between the midpoints of the swings). For every
$N$ one batched run holds all local rates at once
(\IfFileExists{generated/permutational.tex}{$\gamma\in\{\PMGe\}$, $\gamma_\phi\in\{\PMGd\}$, $\PMBatch$ parameter sets
for $N=\PMNs$}{ten parameter sets}); a batch element is dropped once its lifetime is settled, and the run ends at
$\kappa t=\IfFileExists{generated/permutational.tex}{\PMT}{45}$ (censored entries are marked $>$). Without local
terms only the sector $j=S$ is populated, and that run uses the single block $M_{N+1}$
(\texttt{PermutationalSystem(N, j\_min=N/2)}). The mean-field limit $\tau_{\mathrm{MF}}(\gamma)$ is the same lifetime for
the equations $\dot m_x=\omega m_z+\kappa m_xm_z-\Gamma_2m_x$, $\dot m_y=\kappa m_ym_z-\Gamma_2m_y$, $\dot
m_z=-\omega m_x-\kappa(m_x^2+m_y^2)-\gamma(1+m_z)$, $\Gamma_2=\gamma/2+2\gamma_\phi$ (drive about $y$).

\begin{table}[ht]
\centering\small
\caption{Lifetime $\kappa\tau$ of the oscillations of $m_z$ (parameters of the text): without local terms ($0$), with
local emission $\gamma/\kappa$ (columns 3--7) and with local dephasing $\gamma_\phi/\kappa$ (columns 8--12); last row:
mean field ($N\to\infty$). $>$: no lifetime before the end of the run.}\label{tab:btc}
\IfFileExists{generated/permutational.tex}{%
\begin{tabular}{r|c|ccccc|ccccc}
\toprule
\PMTauHeader\\
\midrule
\PMTauRows
\bottomrule
\end{tabular}}{\fbox{\small generated/permutational.tex not found}}
\end{table}

\begin{figure}[ht]
\centering
\fig{permutational.pdf}{\textwidth}
\caption{(a) $m_z(t)$ for the largest $N$ without and with local emission, and the mean-field curves (dotted).
(b) $\tau(N,\gamma)$ for local emission; dotted: $\tau_{\mathrm{MF}}(\gamma)$, dashed: the additive
rule~\eqref{eq:btc-rule}. (c) All lifetimes with a finite mean-field limit, $\tau/\tau_{\mathrm{MF}}$ against
$x=\tau(N,0)/\tau_{\mathrm{MF}}$; the dashed curve $x/(1+x)$ is~\eqref{eq:btc-rule}. (d) The relative deviation from
the rule; points of equal rate are joined, $N$ increasing to the right (filled: emission, open: dephasing).}\label{fig:btc}
\end{figure}

\paragraph{Results.} \IfFileExists{generated/permutational.tex}{%
Table~\ref{tab:btc} and Figure~\ref{fig:btc}. Without local terms the lifetime grows linearly,
$\kappa\tau(N,0)\approx\PMSlope\,N+\PMIntercept$ for $\PMFitNmin\le N\le\PMNmax$ (least squares). The growth is
linear, as expected: for $\omega\gg\kappa$ the oscillations of $m_z$ from all spins down decay at the rate
$3\kappa/(2N)$~\cite{liu2026,nemeth2025}, which corresponds to $\kappa\tau\approx2N/3$; our $\omega=\PMOmega\kappa$ lies
outside that perturbative regime, and the slope is smaller. With local
terms it saturates at the mean-field value (last row of Table~\ref{tab:btc}), and the two damping mechanisms combine
approximately like independent decay rates,
\begin{equation}\label{eq:btc-rule}
  \frac1{\tau(N,\gamma)}\approx\frac1{\tau_{\mathrm{rule}}}=\frac1{\tau(N,0)}+\frac1{\tau_{\mathrm{MF}}(\gamma)} ,
\end{equation}
for local emission and local dephasing alike: in the variables $x=\tau(N,0)/\tau_{\mathrm{MF}}$ and
$\tau/\tau_{\mathrm{MF}}$ all data lie close to the curve $x/(1+x)$ (Figure~\ref{fig:btc}c), and the crossover size at
which the local damping takes over ($x=1$) is $N^*\approx\kappa\tau_{\mathrm{MF}}(\gamma)/\PMSlope$. Over the
$\PMRuleCount$ entries with finite $\tau$ and $\tau_{\mathrm{MF}}$, $\tau/\tau_{\mathrm{rule}}-1$ lies between
$\PMRuleMin$ and $\PMRuleMax$, and between $\PMRuleBigMin$ and $\PMRuleBigMax$ for $N\ge\PMRuleNmin$. The deviation has a
structure (Figure~\ref{fig:btc}d): for the larger rates $\tau$ exceeds $\tau_{\mathrm{rule}}$, by an amount that
decreases at large $N$; for the smallest rates the deviation decreases through zero as $N$ grows, so that there
$\tau<\tau_{\mathrm{rule}}$ at the largest $N$; at equal $x$ local dephasing deviates more than local emission.
So~\eqref{eq:btc-rule} describes the lifetimes to within the range just stated; it is not an asymptotic law that
our data establish.%
}{Run \texttt{scripts/permutational.py}.}

\paragraph{What is reproduced and what is new.} Without local terms, the growth of the lifetime with $N$ is known:
Iemini et al.~\cite{iemini2018} report a decay rate that vanishes as a power of $N$, Piccitto
et al.~\cite{piccitto2021} state that it decreases as $N^{-1}$ for this model, and for $\omega\gg\kappa$ Liu
et al.~\cite{liu2026} and Nemeth et al.~\cite{nemeth2025} derive the rate $3\kappa/(2N)$; only the slope
\IfFileExists{generated/permutational.tex}{\PMSlope}{0.45} at
$\omega=\IfFileExists{generated/permutational.tex}{\PMOmega}{2}\kappa$ is, as far as we know, new. The damping of the
mean-field limit cycle by non-collective dissipation was found by Piccitto et al.~\cite{piccitto2021} (for local
decay see also~\cite{wang2025}), and for local terms it follows directly from the mean-field equations above; that
local dephasing and local losses wash out the oscillations at finite $N$ was shown with PIQS for
$N=30$~\cite{shammah2018}. The superradiance checks reproduce textbook results~\cite{dicke1954,gross1982}. The
finite-$N$ lifetimes with local
emission or dephasing for $N$ up to \IfFileExists{generated/permutational.tex}{\PMNmax}{80} (Table~\ref{tab:btc}), the
approximate additivity~\eqref{eq:btc-rule} of the two damping rates and the structure of its deviation are, as far as
we know, new numbers. They are observations for this model, these rates and this definition of the lifetime; we have
no proof of~\eqref{eq:btc-rule}.

\paragraph{Cost.} \IfFileExists{generated/permutational.tex}{%
The batched run for $N=\PMNmax$ ($\PMBatch$ parameter sets) took $\PMSecNmax$~s, the whole sweep $\PMSecSweep$~s and
the script with all checks $\PMSeconds$~s on one CPU thread (\PMMode\ mode, \PMDevice; the full mode, with $N$ up to
$200$, is meant for a GPU, the cost per step growing as $N^3$ and the number of steps as $N$). For $N=\PMPiqsNmax$
one trajectory of the validation took $\PMPiqsSecOurs$~s here and $\PMPiqsSecBuild+\PMPiqsSecSolve$~s with PIQS
(Liouvillian, \texttt{mesolve}) at matched accuracy, and the generic route with the dense
\texttt{InterSectorChannel} needs $2C^2k_{\max}^2$ Kraus entries ($C=N/2+1$ sectors, $k_{\max}=N+1$), which limits it
to a few tens of spins. On $(\C^2)^{\otimes N}$ the state alone would have $\PMDimFullNmax$ entries for
$N=\PMNmax$.%
}{}
